\begin{statementbox}
设 $I$ 为 $\triangle ABC$ 的内心，角 $A, B, C$ 所对的边长分别为 $a, b, c$，则有以下核心结论：
\begin{enumerate}
    \item \textbf{基准等式}：$a\vec{IA} + b\vec{IB} + c\vec{IC} = \vec{0}$；
    \item \textbf{任意点表示}：对于平面内任意一点 $O$，有 $\vec{OI} = \frac{a\vec{OA} + b\vec{OB} + c\vec{OC}}{a+b+c}$；
    \item \textbf{角平分线向量}：若 $AD$ 为 $\angle A$ 的平分线，则 $\vec{AD} = \frac{b}{b+c}\vec{AB} + \frac{c}{b+c}\vec{AC}$；
    \item \textbf{单位向量表示}：向量 $\vec{v} = \lambda \left( \frac{\vec{AB}}{|\vec{AB}|} + \frac{\vec{AC}}{|\vec{AC}|} \right) (\lambda \neq 0)$ 所在直线过 $\triangle ABC$ 的内心 $I$（即该向量与角平分线 $AD$ 共线）。
    \item \textbf{内心坐标公式}：在平面直角坐标系中，设 $A(x_1, y_1), B(x_2, y_2), C(x_3, y_3)$，则内心 $I$ 的坐标为：
    \[ I\left( \frac{ax_1 + bx_2 + cx_3}{a+b+c}, \frac{ay_1 + by_2 + cy_3}{a+b+c} \right) \]
\end{enumerate}
\tcbline % 如果你使用的是 tcolorbox，可以用这条分割线

\textbf{延伸思考：}
\begin{itemize}
    \item \textbf{减法形式}：若将加号改为减号，即向量 $\vec{u} = \lambda \left( \frac{\vec{AB}}{|\vec{AB}|} - \frac{\vec{AC}}{|\vec{AC}|} \right) (\lambda \neq 0)$，则该向量所在直线平分 $\triangle ABC$ 的\textbf{外角}。
    \item \textbf{几何直观}：加法形式构造的是菱形的对角线（内角平分线），减法形式构造的是与该对角线垂直的另一条对角线（外角平分线）。
\end{itemize}
\end{statementbox}